\section{Method}

\subsection{Setting}

At each communication round, the server broadcasts the current global model to
all clients. Anchor clients return full-precision updates, while constrained
clients return low-bit quantized updates. The server's job is to repair the
constrained uploads before aggregation without disturbing the clean anchor
signal.

\begin{figure*}[t]
  \centering
  {\def\methoddiagramfont{\normalsize}%
  \resizebox{0.99\textwidth}{!}{\providecommand{\methoddiagramfont}{\small}

\begin{tikzpicture}[
  >=Stealth,
  font=\methoddiagramfont,
  box/.style={
    draw,
    rounded corners=2.5pt,
    align=center,
    inner sep=4pt,
    minimum height=9mm
  },
  model/.style={box, fill=black!5, draw=black!45, text width=2.5cm},
  anchorbox/.style={box, fill=anchorblue!12, draw=anchorblue, text width=2.9cm},
  quantbox/.style={box, fill=quantred!12, draw=quantred, text width=2.9cm},
  serverbox/.style={box, fill=servergold!12, draw=servergold!80!black, text width=3.0cm},
  bufferbox/.style={box, fill=buffergreen!12, draw=buffergreen!70!black, text width=3.0cm}
]

\node[model] (global) at (0, 0) {Global model};

\node[anchorbox] (anchors) at (-5.4, -2.3) {Anchor clients\\full-precision uploads};
\node[serverbox] (targets) at (0, -2.3) {Target matching\\mean anchor or nearest anchor};
\node[quantbox] (quantized) at (5.4, -2.3) {Constrained clients\\low-bit uploads};

\node[bufferbox] (replay) at (-3.8, -4.9) {Replay buffer\\keep recent supervision pairs};
\node[serverbox] (corrnet) at (0, -4.9) {CorrectionNet\\residual MLP learns correction};
\node[serverbox] (blend) at (3.8, -4.9) {Blend with original update\\before aggregation};

\node[serverbox] (agg) at (0, -7.7) {Aggregate anchors and corrected uploads};
\node[model] (next) at (5.5, -7.7) {Updated model};

\coordinate (aggleft) at ([xshift=-0.8cm]agg.west);
\coordinate (aggright) at ([xshift=0.8cm]agg.east);
\coordinate (blenddrop) at (blend.south |- aggright);

\draw[->, thick] (global.south west) -- (anchors.north);
\draw[->, thick] (global.south east) -- (quantized.north);
\draw[->, thick, anchorblue] (anchors.east) -- node[above, font=\footnotesize] {supervision} (targets.west);
\draw[->, thick, quantred] (quantized.west) -- node[above, font=\footnotesize] {low-bit uploads} (targets.east);
\draw[->, thick] (targets.south west) -- (replay.north east);
\draw[->, thick] (replay.east) -- (corrnet.west);
\draw[->, thick] (corrnet.east) -- (blend.west);
\draw[->, thick] (blend.south) -- node[pos=0.45, right, font=\footnotesize] {corrected uploads} (blenddrop);
\draw[->, thick] (blenddrop) -- (aggright) -- (agg.east);
\draw[->, thick, anchorblue, dashed] (anchors.south) |- node[pos=0.74, below, font=\footnotesize] {anchor uploads} (aggleft) -- (agg.west);
\draw[->, thick] (agg.east) -- (next.west);

\end{tikzpicture}
}}
  \caption{Implemented round-level pipeline for anchor-guided dequantization.
  Full-precision anchor clients supervise a server-side correction model. The
  server matches constrained uploads to anchor targets, reuses recent
  supervision through a replay buffer, applies a learned correction, blends it
  with the original low-bit update, and aggregates the repaired uploads with
  the anchors.}
  \label{fig:method_architecture}
\end{figure*}

\subsection{Target Construction}

The correction module operates only on \emph{learnable} parameters. BatchNorm
running statistics bypass the correction path unchanged. Each client delta is
flattened into a single vector before training the corrector.

Two target modes are implemented. In the simpler version, every constrained
client is trained toward the mean anchor update for that round. In the smart
matching version, each constrained client is paired with the anchor whose update
has the highest cosine similarity. The smart-matching version is the one
reported as \method.

\subsection{Correction Network and Replay Buffer}

The server trains a residual MLP, CorrectionNet, on pairs of quantized uploads
and anchor-derived targets. The network uses one encoder layer, three residual
blocks, and a zero-initialized decoder so that it starts close to the identity
map and only has to learn a correction rather than a full replacement update.

Instead of training only on the current round, the server appends each new
mini-batch of supervision pairs to a replay buffer capped at the most recent 10
rounds. All buffered pairs are concatenated for training. Before the correction
update, buffered quantized deltas are normalized with a single scalar
$z$-score computed across the buffer so that the network sees a more stable
input scale from round to round.

\subsection{Training Objective}

The base training signal is a reconstruction loss that pushes corrected uploads
toward their anchor-derived targets. When a server validation loader is
available, the implementation can also add a distillation term. In that case,
the current global model acts as the teacher, and the corrected update is also
encouraged to produce similar logits on the validation split.

The teacher logits come from the current global model on the validation split.
The student logits come from a model built by applying the mean corrected delta
from the current round. The reported configurations use distillation
temperature $\tau=4$.

\subsection{Applying Correction and Aggregation}

At inference time, the server denormalizes the network output and can blend the
corrected upload with the original quantized upload. The main smart-matching
configuration uses an 80/20 mix of corrected and original update, which helps
preserve client-specific direction under non-IID data.

The corrected deltas are then aggregated together with anchor deltas. The code
supports either plain sample-count weighting or a convex combination of
sample-count and cosine-similarity weights relative to the mean anchor delta.
Finally, the aggregated update is applied to the global model.

\subsection{Warmup}

The smart-matching strategy also uses a warmup phase. For the first five rounds,
the server trains the CorrectionNet and fills the replay buffer but still
aggregates the raw client deltas with FedAvg. Correction is only applied after
the buffer contains enough supervision signal to avoid harming early
optimization.
